\section{Transferencia Markov compatible con todo el límite inverso}
\label{sec:cpe-transferencia-limite-inverso}

Sea \(\rho_m:H_m\twoheadrightarrow Y_m\) el cociente observable a profundidad
\(m\) y sea \(L_m(y,y')\) el núcleo de transición ya construido en
\(Y_m\). Para cada fibra finita \(\rho_m^{-1}(y)\), sea
\(\nu_{m,y}\) la medida uniforme condicionada. La regularidad nonádica de
las extensiones permite escoger estas medidas de modo que su empuje por el
truncamiento sea \(\nu_{m-1,\bar y}\).

Defínase, para \(h\in H_m\),
\begin{equation}
 \mathcal K_m(h,A)
 :=\sum_{y'\in Y_m}
 L_m(\rho_m(h),y')\,
 \nu_{m,y'}\bigl(A\cap\rho_m^{-1}(y')\bigr).
 \label{eq:cpe-kernel-lift-uniforme}
\end{equation}

\begin{theorem}[Elevación fuertemente agregable y límite de Markov]
\label{thm:cpe-kernel-proyectivo}
Los núcleos \(\mathcal K_m\) satisfacen
\begin{equation}
 \sum_{\rho_m(g)=y'}\mathcal K_m(h,\{g\})
 =L_m(\rho_m(h),y'),
 \label{eq:cpe-agregabilidad-fuerte}
\end{equation}
independientemente del representante \(h\) de la fibra. Además, conmutan con
los truncamientos. Existe, por tanto, un único núcleo de Markov
\(\mathcal K_\infty\) sobre el límite inverso cuya marginal de nivel \(m\)
es \(\mathcal K_m\). Su proyección por \(\rho_\infty\) es el núcleo
observable \(L_\infty\).
\end{theorem}

\begin{proof}
La suma del miembro izquierdo de
\eqref{eq:cpe-agregabilidad-fuerte} es
\[
 L_m(\rho_m(h),y')
 \sum_{g\in\rho_m^{-1}(y')}\nu_{m,y'}(\{g\})
 =L_m(\rho_m(h),y').
\]
La compatibilidad de \(L_m\), de \(\rho_m\) y de las medidas condicionadas
prueba el cuadrado de truncación. Las distribuciones cilíndricas obtenidas
iterando \(\mathcal K_m\) son consistentes; el teorema de extensión para
sistemas proyectivos de espacios finitos produce \(\mathcal K_\infty\), y
la unicidad sigue de que los cilindros generan la sigma-álgebra del límite.
\end{proof}

La construcción conserva la quietud microscópica y las etiquetas completas
en la medida condicionada. El promedio visible no reemplaza el estado
enriquecido y no se identifica con el Topological Phase Kernel.

